\section{Preliminaries and Setting}
\label{sec:prelim}

\subsection{Words as bit-polynomials}

Fix a word width $\wbit$ (deployed: $\wbit = 32$) and the \emph{cell ring}
\[
  \cellring \;:=\; \Ftwo[X]/(X^{\wbit}),
\]
a free $\Ftwo$-module of rank $\wbit$. We identify $\cellring$ with the
set of bit-polynomials $\bitpoly{\wbit}$: a $\wbit$-bit word
$u = (u_0,\dots,u_{\wbit-1})$ is the polynomial $\bp{u} = \sum_b u_b X^b$
whose coefficient of $X^b$ is the bit $u_b$. Coefficient-wise addition in
$\cellring$ is bitwise $\XOR$ (so $1+1=0$), and the constant $1$ is the
all-zero-but-LSB word.

Two further quotients capture the word-level operations that move bits
between positions:
\[
\begin{gathered}
  \Fshift \;=\; \Ftwo[X]/(X^{\wbit})
  \quad\text{(shift ring),}\\
  \Frot \;=\; \Ftwo[X]/(X^{\wbit}-1)
  \quad\text{(rotation ring).}
\end{gathered}
\]
In $\Fshift$, multiplication by $X^c$ is a logical left shift by $c$ with
zero-fill (high bits fall off the top); in $\Frot$, multiplication by
$X^c$ is a cyclic rotation. The cell ring $\cellring$ \emph{is} the shift
ring; we use the separate name $\Fshift$ when the shift structure is the
point. These three rings share the underlying $\Ftwo$-module
$\bitpoly{\wbit}$ and differ only in how multiplication wraps.

\subsection{The binary extension field}

Fix a binary extension field $\K = \mathrm{GF}(2^{\kappa})$ (deployed:
$\kappa = 128$, with defining pentanomial $X^{128}+X^7+X^2+X+1$). $\K$ is
the protocol's \emph{challenge and evaluation field}: all Fiat--Shamir
randomness and all sumcheck arithmetic live in $\K$. Addition in $\K$ is
$\XOR$ of $\kappa$-bit vectors; multiplication is a carryless product
reduced modulo the defining polynomial. $\K$ contains $\Ftwo$ as its
prime field, and we use the trivial embedding $\Ftwo \hookrightarrow \K$
throughout (a bit $b\in\{0,1\}$ is the field element $b\cdot 1_\K$).

\begin{remark}[The two roles of $X$]
\label{rem:two-X}
The symbol $X$ is the \emph{cell variable}: it indexes the $\wbit$ bit
positions of a word and, once coefficients are lifted into $\K$, the
$\wbit$ coordinates of a polynomial in $\Kpoly$. It is distinct from the
formal variable defining $\K$. Every projection below substitutes the
cell variable $X$ by an element of $\K$; none touches $\K$'s internal
representation. We keep the single letter $X$ and rely on the ring
context.
\end{remark}

\subsection{The trace model}
\label{sec:trace}

The computation is presented as an algebraic intermediate representation
(AIR) with lookups --- in the codebase, a \emph{universal AIR} (UAIR).
A trace is a matrix of cells
\[
  T \;\in\; \cellring^{\,N_{\mathrm{rows}} \times N_{\mathrm{cols}}},
\]
with $N_{\mathrm{rows}} = 2^{\nu}$ rows (the \emph{steps} of the
computation) and $N_{\mathrm{cols}}$ columns. Each column
$\col{c}\in\cellring^{\,N_{\mathrm{rows}}}$ is a vector of words, one per
step; we write $\col{c}[i]$ for its value at row $i$ and $\col{c}[i\!+\!1]$
for the next row, so transition constraints relate consecutive rows.
Columns are partitioned into:
\begin{itemize}[leftmargin=2em]
  \item \emph{public} columns (round constants, padding, the message),
    known to the verifier and not committed;
  \item \emph{primary witness} columns, committed by the prover; and
  \item \emph{virtual} columns, fixed $\Ftwo$-linear functions of a
    primary column --- shifts, rotations, and $\XOR$-combinations of these
    (such as the $\Sigma/\sigma$ mixing maps) --- reconstructed rather than
    committed (\Cref{sec:arith}).
\end{itemize}

\paragraph{Constraints.} A UAIR imposes two kinds of constraint, each
required to hold at every row (transition constraints additionally
reference row $i\!+\!1$):
\begin{itemize}[leftmargin=2em]
  \item \emph{Linear / ideal-membership constraints.} A constraint is a
    relation $Q(T[i],T[i\!+\!1]) \in \ideal{g}$ asserting that an
    $\Ftwo[X]$-combination of cells lies in the ideal generated by $g$.
    Equality is the case $g = 0$ (e.g.\ $\col{c} = \col{a}+\col{b}$ in
    $\cellring$); the modular-addition carry contributes the bit-$0$ case
    $g = X$; and word width is enforced by working modulo $g = X^{\wbit}$.
    Rotations and shifts do not appear here --- they are realised as
    virtual columns (above), not constraints. These are the constraints
    that survive the linear machinery untouched.
  \item \emph{Hadamard (Boolean-product) constraints.} A relation
    $\col{c} = \col{a}\had\col{b}$, where $\had$ is coefficient-wise
    multiplication of bit-polynomials --- i.e.\ bitwise $\AND$. This is
    the only nonlinear constraint type, and the only one that does not
    pass through the commitment linearly.
\end{itemize}
A trace \emph{satisfies} the UAIR if all constraints hold at all rows.
The arithmetization of \Cref{sec:arith} expresses SHA-256 in exactly this
language.

\subsection{Multilinear extensions}

We index rows by the boolean hypercube $\{0,1\}^{\nu}$. For a column
$v\in\cellring^{\,N_{\mathrm{rows}}}$ and any $\Ftwo$-algebra $S$
containing $\cellring$, the \emph{multilinear extension} (MLE) of $v$
over $S$ is the unique multilinear $\mle{v}\in S[Y_1,\dots,Y_\nu]$ with
$\mle{v}(c) = v_c$ for all $c\in\{0,1\}^\nu$; explicitly
\[
  \mle{v}(r) \;=\; \sum_{c\in\{0,1\}^{\nu}} \eq_c(r)\,v_c,
  \qquad
  \eq_c(r) = \prod_{k=1}^{\nu}\bigl(c_k r_k + (1-c_k)(1-r_k)\bigr).
\]
The PIOP works with MLEs whose coefficients are projected into $\K$ (so
$S = \K$) and, transiently, with MLEs valued in $\cellring$ or $\Kpoly$.

\subsection{The cell projection: one evaluation, two bases}
\label{sec:projections}

The reduction from $\Ftwo[X]$-typed constraints to a field collapses each
cell to a single element of $\K$ by an $\Ftwo$-linear functional. The
protocol samples only \emph{one} such functional --- evaluation at a random
point $\alpha$ --- but reads the witness in two bases, as made precise
below. We state the general functional first.

\begin{definition}[Cell projection]
\label{def:projection}
A \emph{cell projection} is an $\Ftwo$-linear map
$\psi_\xi:\cellring\to\K$ determined by a weight vector
$\xi = (\xi_0,\dots,\xi_{\wbit-1})\in\K^{\wbit}$ via
\[
  \psi_\xi\Big(\textstyle\sum_b w_b X^b\Big) \;=\; \sum_b w_b\,\xi_b.
\]
It extends $\K$-linearly to $\Kpoly$ by $\psi_\xi(\sum_b p_b X^b) =
\sum_b p_b\,\xi_b$ for $p_b\in\K$, i.e.\ the inner product of the
coefficient vector with $\xi$.
\end{definition}

The system uses a single evaluation point: both functionals it needs are
the substitution $X=\alpha$ for one transcript-fresh $\alpha\in\K$, applied
to the witness read in two different bases.
\begin{itemize}[leftmargin=2em]
  \item The \emph{monomial reading}
    $\psia(w) = \bigl(\textstyle\sum_b w_b X^b\bigr)\big|_{X=\alpha} = \sum_b w_b\alpha^b$
    is genuine evaluation $w(\alpha)$, an $\Ftwo$-\emph{algebra} homomorphism
    $\cellring\to\K$: it respects multiplication by $X$ --- i.e.\ shifts ---
    and the cell ideals, and so discharges the \emph{linear} constraints
    (\Cref{sec:piop}).
  \item The \emph{Lagrange reading}
    $\psiz(w) = \bigl(\textstyle\sum_b w_b L_b(X)\bigr)\big|_{X=\alpha} = \sum_b w_b L_b(\alpha)$,
    where $\{L_b\}$ is the Lagrange basis of an $\Ftwo$-linear evaluation
    domain $H\subset\K$ of size $\wbit$, turns the coefficient-wise
    (Hadamard) product into a pointwise product on $H$, and so discharges
    the one nonlinear constraint, $\AND$ (\Cref{sec:hadamard}).
\end{itemize}
These are \emph{not} two projections. There is one substitution $X=\alpha$;
$\psia$ and $\psiz$ are its values when the bits are laid into the monomial
resp.\ Lagrange basis, so even the weight vectors coincide as
$\alpha^b = X^b|_{\alpha}$ and $L_b(\alpha) = L_b(X)|_{\alpha}$ --- the same
point in two bases. The basis is \emph{forced by the constraint} --- monomial
where multiplication-by-$X$ structure matters, Lagrange where a
coefficient-wise product does --- not chosen freely. Because the commitment
of \Cref{sec:commitment} binds the bit vector itself (its bit-slice fold),
each reading is one length-$\wbit$ inner product against that one bound
object, at the one point $\alpha$ --- which is what lets the linear and
Hadamard parts share a single opening.

\begin{remark}[Why $\Ftwo$-linearity is the through-line]
\label{rem:linearity-thread}
Each $\psi_\xi$ is $\Ftwo$-linear, and the linear code used by the
commitment is $\Ftwo$-linear and does not widen $\cellring$. Hence
projections commute through encoding, $\XOR$ of columns commutes with
projection, and shift/rotation virtual columns commute with both. This
single fact is reused at every layer: it makes $\XOR$ free, makes shifts
virtual, lets the random projection be applied after committing, and lets
one commitment serve many functionals.
\end{remark}
